\section{Introducción}

\subsection{Contexto}
Durante el cursado de la asignatura Sistemas III, el equipo desarrolló un sistema de planificación de recursos empresariales (ERP) orientado a la gestión integral de un supermercado.
Dicho sistema abarca módulos de ventas presenciales, cajas, inventario, depósitos y transferencias, compras, proveedores, comercio electrónico y un tablero gerencial con datos históricos sintéticos.
En la modalidad de comercio electrónico, cada compra realizada por un cliente genera un pedido online que debe ser armado dentro del establecimiento antes de quedar disponible para su retiro o despacho.
Este trabajo toma ese entorno operativo como caso de estudio para analizar la etapa de preparación de pedidos online desde la perspectiva de la Investigación Operativa.

\subsection{Problema}
La llegada de pedidos online a lo largo del día presenta variaciones horarias que dificultan dimensionar de manera intuitiva la cantidad de personal necesaria en el sector de preparación.
Cuando la tasa de arribo de pedidos supera la capacidad de atención de los empleados disponibles, se acumulan órdenes pendientes en espera y se incrementa el tiempo total hasta que cada pedido queda listo.
Por el contrario, asignar una dotación excesiva reduce las demoras operativas, pero genera capacidad ociosa durante las franjas de menor actividad e incrementa los costos laborales.
Ante este compromiso entre nivel de servicio y utilización de recursos, el problema central consiste en determinar cuántos preparadores conviene asignar para atender distintos niveles de demanda y cumplir un objetivo de tiempo de preparación definido.

\subsection{Sistema analizado}
El sistema bajo estudio comprende el recorrido operativo de cada pedido online desde su ingreso hasta su finalización en el sector de armado.
El flujo considerado sigue cuatro etapas consecutivas: la llegada del pedido al sistema, la espera en una fila común si todo el personal se encuentra ocupado, la preparación por parte de un empleado disponible y la disponibilidad final del pedido preparado.
Los pedidos se atienden bajo una disciplina de orden de llegada y cada preparador procesa un único pedido a la vez.
El alcance del modelo se circunscribe hasta el instante en que el pedido queda listo, asumiendo disponibilidad de productos en inventario y empleados con desempeño equivalente, por lo que el ruteo interno en el depósito, la gestión de faltantes y el reparto domiciliario quedan fuera del límite de este estudio inicial.

\subsection{Justificación de las técnicas}
Abordar la planificación de personal en este proceso requiere articular tres herramientas complementarias de la Investigación Operativa: pronósticos de demanda, teoría de colas y simulación de eventos discretos.
En primer lugar, los métodos de pronóstico permiten estimar el perfil horario de arribo de pedidos para la jornada siguiente a partir del registro histórico y cuantificar el error de predicción sobre un conjunto de prueba reservado \cite{hyndman_athanasopoulos_fpp3}.
En segundo lugar, la teoría de colas proporciona el marco formal para relacionar la tasa de llegadas, la capacidad de servicio y los tiempos de espera, ofreciendo además el modelo analítico $M/M/c$ como referencia exacta bajo condiciones de régimen estacionario.
Finalmente, la simulación de eventos discretos mediante recursos compartidos de capacidad limitada permite representar perfiles horarios de demanda variable, conservar los pedidos pendientes entre horas sucesivas y comparar el desempeño de distintas dotaciones en escenarios de demanda baja, media y alta \cite{simpy_shared_resources}.
